\section{History and Definitions}
\label{sec:background}
\subsection{Settlement Schemes}
Before the independence of Kenya in 1963, the colonial government implemented a land tenure system favoring European settlers, which led to the displacement of native Kenyans from their ancestral lands.
Especially the fertile highlands today known as \textit{Rift Valley} or \textit{Kenya Highlands} then known as \textit{White Highlands} were reserved for European settlers.
This displacement into \textit{Native Reserves} created a culture of landlessness and poverty among native Kenyans, who traditionally had a system ensuring land ownership to everyone \parencite{leoLandClassKenya1984}.
To address this issue of landlessness and promote agricultural development, the post-colonial government initiated a series of settlement schemes, which aimed to redistribute formerly European-owned land to landless Kenyans \parencite{lukaloExploringStatusSettlement2016}.
Settlement schemes were designed in two main forms: In the beginning, \textit{conventional schemes}, which were located in the former White Highlands and were intended for large-scale commercial farming, were implemented.
One notable example for this category is the Million Acres Scheme, which resettled 35,000 families on 1.2 million acres of land.
Their size and agricultural potential caused a price that only a few wealthy individuals could afford.
After some years passed, and it became clear a large portion of Kenyans could not afford this type of land and with the rebellion against British land politics in mind, \textit{squatter schemes} were implemented.
Squatter schemes, which were much cheaper and located in other, less fertile regions such as the former Native Reserves, were intended for small-scale subsistence farming, and in recent times mostly used to legalize already existing informal settlements \parencite{leoLandClassKenya1984}.
Conventional schemes were further divided into two categories: Low density schemes, which allocated large plots of land to a few settlers, and high density schemes, which allocated smaller plots \parencite{lukaloExploringStatusSettlement2016}.
Allocation of conventional schemes was mostly done in the 1960s, while squatter schemes were the primary focus of land redistribution efforts from 1990 onwards \parencite{lukaloExploringStatusSettlement2016}.
This change in focus can also be attributed to the democratization of Kenya which led to a new constitution in 2010, calling for \textquote{equitable access to land} and a \textquote{security of land rights} \parencite{constitution}.


\begin{customfigure}[t!]
    \caption{Pre-Independent Kenya with the Location of implemented Settlement Schemes}
    \includegraphics[width=0.8\textwidth]{figureDescriptive_Georef}
    \label{fig:desc_map}
    
    \fignotes{This map shows the location of settlement schemes in Kenya (red), overlaid on a map of Kenya published in 1957 by Survey of Kenya.
    Yellow areas represent the former White Highlands, while green areas represent the former \textit{Crown Forest}.}
\end{customfigure}

Still, the issue of land tenure security (LTS) persists, especially in squatter schemes, where only 61\% of the land is registered, compared to 91\% in conventional schemes \parencite{lukaloExploringStatusSettlement2016}.

\subsection{Land Tenure Security}
Land tenure describes the legal and institutional arrangements that govern land ownership and use, including the rights and responsibilities of landholders, for example mechanisms for acquiring and transferring land \parencite{naughton-trevesLandTenureTropical2014}.
This includes various forms of land ownership, such as private ownership, communal ownership, and state ownership \parencite{robinsonDoesSecureLand2014}.
Land tenure can be further broken down into formal \textit{de jure} and informal \textit{de facto} tenure.
De jure refers to state recognized and legally protected land rights often confirmed by the issuance of a land title, while the de facto refers to customary norms and practices, which can be just as binding as formal recognition \parencite{robinsonDoesSecureLand2014, naughton-trevesLandTenureTropical2014}.
The lack of enforcement by some states led to the separation of form and security of land tenure; \cite{naughton-trevesLandTenureTropical2014} defines land tenure security as the degree to which law and society uphold those property rights.
\cite{tsengInfluenceLandTenure2020} defines land tenure security as the alignment of de facto and de jure land tenure.

In addition to the distinction between formal and informal systems, \cite{robinsonDoesSecureLand2014} also differentiate between legitimate and perceived security.
Applying these theoretical concepts to the Rift Valley region reveals how de jure and de facto land rights have been historically unclear and used as a political tool.
Pre-colonial land tenure was governed by customary norms, with patriarchal figures adjusting plot sizes based on the needs of the individual \parencite{leoLandClassKenya1984}.
During the colonial period, the British administration restricted de jure ownership in the Rift Valley to European settlers, with native Kenyans only allowed to cultivate land in this region while working for the settlers \parencite{leoLandClassKenya1984}.
This uncertainty of land tenure continued after independence, as land allocated in settlement schemes was (as of 2016), on average, registered 17 years after allocation with over 60 settlement schemes still having no registration \parencite{boonePromisedLandSettlement2021}.
Without registration, the landholders have no legal claim to their land, which can lead to potential conflicts over land ownership, a higher perceived risk of eviction, and a lack of access to credit.
In Kenya's period of dictatorship under President Daniel arap Moi, this allocation without registration was used as a mechanism for political patronage by revoking land rights of political opponents \parencite{booneLandConflictDistributive2012}.
For many Kenyans, the history of land tenure insecurity has created a sense of distrust towards the government and its ability to protect their land rights \parencite{placeEconomicEffectsLand1998}.

\subsection{Deforestation and Forest Degradation}
Historically, deforestation and forest degradation have been a significant issue in Kenya.
Before colonial rule, forests were managed by local communities and in the region of the Mau Forest, large-scale agriculture was not practiced \parencite{leoLandClassKenya1984}.
As mentioned, the colonial government implemented a land tenure system that distributed large parts of fertile land to European settlers for commercial agriculture.
The establishment of large-scale farming operations led to the clearing of forests which contributed to the loss of about 900,000 hectares of forest in Kenya between 1900 and 1975 \parencite{gitauColonialeraSettlementsPostcolonial2025}.
To address the issue, the colonial government implemented multiple systems with the goal of regenerating tropical forests.
One notable example, the Shamba system, granted native Kenyans access to farming land for a period of a few years in exchange for planting trees.
Overall those systems of exploitation of the forests by European settlers and subsequent regeneration by native Kenyans contributed to the systematic distrust of the local population towards government-mandated forest management \parencite{gitauColonialeraSettlementsPostcolonial2025}.

\begin{customtable}
    \caption{Forest Cover Changes in Kenya's Tropical Moist Forest for different Subgroups since 1990}
    \label{tab:desc_forest_stats}
    \tablerendering{tableForests_Descriptive}
    \tabnotes{This table shows the total number (in thousands) and percentage (relative to all of Kenya) of deforested and degraded pixels in different subgroups using the EU JRC TMF dataset.
    The first row shows the total number of deforested and degraded pixels in Kenya.
    The second row displays observations inside the area used for the study, which includes the settlement schemes and a 3km buffer around them.
    The third row only includes pixels inside the settlement schemes.
    In the fourth row, the total numbers of the first row are shown relative to the forest cover size in 1990, which is the baseline year of the EU JRC TMF dataset.}
\end{customtable}

But not only the historic view on deforestation and overexploitative land use is important, but also recent changes.
Tropical moist forests (TMF) like the Mau Forest in Kenya are one of the most important ecosystems on Earth.
In 2025 about 1.19\% of the Kenya's total land area was covered by tropical moist forests, compared to 2.06\% in 1990.
As can be seen in \cref{tab:desc_forest_stats}, in the time between 1990 and 2025, about 34\% of the total forest area has been deforested, while 39\% has been degraded.
This is not unusual, as many countries have experienced similar forest cover losses.
However, at the same time, the settlement schemes and their surrounding areas have experienced a much higher rate of deforestation and degradation,
with 15\% of deforestation and 20\% of degradation occurring in the study area, an area that only represents 2\% of the total land area of Kenya, according to \cref{tab:desc_forest_stats}.